\section{Problema de ingeniería y selección del método}
\label{sec:introduccion}

La caracterización del terreno responde a dos necesidades distintas de una obra civil. Antes de construir, conocer la variación de resistencia y rigidez con la profundidad es un requisito para dimensionar la fundación y evaluar asentamientos. Durante la construcción y una vez terminada la obra, vuelve a ser necesario verificar que el suelo, los rellenos o las capas compactadas conserven la uniformidad prevista y detectar cambios que justifiquen mantenimiento. El primer problema condiciona el diseño; el segundo exige mediciones repetibles que no dañen lo ya construido \citep{FHWA2006,Park2018Roads}.

Los procedimientos tradicionales establecen contacto directo con el suelo. El ensayo de penetración estándar (SPT) mide resistencia a la penetración y recupera muestras en puntos discretos; CPT y CPTu proporcionan un registro vertical más continuo, aunque pierden aplicabilidad en gravas o estratos cementados. Los ensayos \emph{downhole} y \emph{crosshole} permiten estimar velocidades de propagación entre puntos conocidos, pero requieren uno o varios pozos. Son técnicas consolidadas y necesarias cuando se requieren muestras o medidas localizadas; su costo y carácter puntual dificultan repetirlas sobre extensiones grandes o sobre una obra terminada \citep{FHWA2006,Foti2014}.

Una alternativa no invasiva consiste en excitar el terreno desde la superficie y observar cómo se propagan en él ondas mecánicas. Los tiempos de llegada, las amplitudes y las diferencias de fase contienen información sobre las propiedades elásticas del medio atravesado. Esta familia no reemplaza por principio a los sondeos: aporta una descripción continua y repetible que puede complementarlos y que resulta especialmente útil cuando no se admite perforar o se desea comparar el mismo emplazamiento a lo largo del tiempo \citep{Foti2014,Foti2018}.

Entre los procedimientos de superficie, el análisis multicanal de ondas superficiales (MASW) registra una excitación controlada mediante receptores dispuestos a diferentes distancias. Su interés para este proyecto es doble. Desde la geotecnia, aprovecha la dispersión de las ondas Rayleigh para relacionar frecuencia y rigidez de las capas someras. Desde la instrumentación, exige preservar simultáneamente banda, fase, margen dinámico, referencia temporal, geometría y trazabilidad de los registros. Por ello ofrece un caso de uso completo para diseñar y evaluar una plataforma electrónica de adquisición distribuida \citep{Park1999,Foti2014}.

El objetivo de esta primera etapa es diseñar e implementar esa cadena de medida, desde la fuente instrumentada y el geófono hasta el acondicionamiento, la digitalización, el almacenamiento y la formación de un registro ordenado por distancia. La interpretación mediante inversión se reserva para la siguiente etapa: antes de estimar un modelo del terreno debe demostrarse qué magnitud se capta, cómo se conserva y con qué limitaciones instrumentales.

